\begin{table*}[t]
\centering
\caption{Comparison between related work and this study.}
\label{tab:related_work}
\footnotesize
\setlength{\tabcolsep}{4pt}
\resizebox{\textwidth}{!}{%
\begin{tabular}{p{2.7cm}p{3.7cm}p{4.2cm}p{5.0cm}}
\toprule
\textbf{Work} & \textbf{Main focus} & \textbf{Treatment of participation or cost} & \textbf{Limitation addressed by this study} \\
\midrule

McMahan et al.~\cite{mcmahan2017communication} &
FedAvg and communication-efficient FL. &
Uses partial client participation to reduce communication overhead. &
Does not isolate the \textit{join ratio} as the main variable nor evaluate the accuracy--cost trade-off. \\

Cho et al.~\cite{cho2020client} &
Client selection and convergence in FL. &
Shows that selected clients can affect convergence and model quality. &
Focuses on selection strategies, not on fixed \textit{join ratio} levels and their cost implications. \\

Maciel et al.~\cite{maciel2024federated} &
Energy-aware client selection. &
Uses participant selection to reduce energy consumption in FL. &
Targets adaptive selection, while this study quantifies the effect of predefined \textit{join ratio} values. \\

Qiu et al.~\cite{qiu2023carbon} &
Carbon footprint of FL. &
Evaluates environmental impact associated with federated training. &
Does not focus on the \textit{join ratio} as a design variable for accuracy, time, and energy trade-offs. \\

Savazzi et al.~\cite{savazzi2023energy} &
Energy and carbon analysis of distributed and federated learning. &
Analyzes resource and environmental costs in distributed training. &
Provides a broader energy analysis, but does not isolate fixed participation levels across FL algorithms. \\

\midrule

\textbf{This work} &
\textbf{Impact of client participation on the performance--cost trade-off in FL.} &
\textbf{Treats the \textit{join ratio} as the main experimental variable and evaluates best test accuracy, execution time, estimated energy, and efficiency.} &
\textbf{Shows how increasing participation leads to diminishing returns under a controlled comparison of FedAvg, FedProx, and SCAFFOLD.} \\

\bottomrule
\end{tabular}%
}
\end{table*}
